\chapter[1901 --- Behring]{1901: Emil von Behring}
\label{ch:1901_emilvonbehring}

\section*{Main Contribution}

\textit{Discovered that blood serum from immunized animals could confer passive immunity, leading to the first effective therapy for diphtheria---saving tens of thousands of lives.}

\section{Official Citation}

for his work on serum therapy, especially against diphtheria

\section{Background and Historical Context}

Emil Adolf von Behring was born on March 15, 1854, in Hansdorf, West Prussia (now Laski, Poland), the eldest of thirteen children in a family of modest means\index{Behring, Emil von|textbf}. His father, August Behring, was a schoolmaster who instilled in his children a deep appreciation for education and disciplined study. Young Emil initially studied medicine at the Army Medical School in Berlin, graduating in 1878. His early career was shaped by military medicine; he served in the Prussian army and was stationed at various military hospitals where he encountered the devastating effects of infectious diseases on soldiers.

The late nineteenth century was a period of tremendous upheaval in the understanding of infectious disease. Louis Pasteur's germ theory had firmly established that microorganisms caused disease, and Robert Koch had formulated his famous postulates for proving that a specific organism caused a specific disease\index{Robert Koch|textbf}. Yet the practical challenge of combating infections remained formidable. Vaccines existed for only a handful of diseases, and the therapeutic arsenal available to physicians was shockingly limited.

Before Behring's work, several physicians had attempted to treat infectious diseases using blood transfusions or serum from convalescent patients. In 1890, the French physician Alfred Boinet reported successful treatment of diphtheria using blood from recovered patients, though his results were modest and not widely adopted\index{Alfred Boinet|textbf}. Simultaneously, the Italian physician Camillo Castelli in Buenos Aires was experimenting with antitoxic sera for anthrax and cholera, while the French bacteriologist Jean-Baptiste Perrin demonstrated that heat-inactivated serum could protect against cholera\index{Jean-Baptiste Perrin|textbf}. These parallel efforts suggested that the scientific community was ready for a breakthrough in serum therapy.

Diphtheria, in particular, was one of the most feared childhood illnesses in industrialized nations\index{Diphtheria|textbf}. Caused by the bacterium \emph{Corynebacterium diphtheriae}, the disease produced a thick pseudomembrane in the throat that could asphyxiate victims, and a potent toxin that damaged the heart and nervous system. In the 1880s and 1890s, diphtheria epidemics swept through Europe and North America, killing tens of thousands of children annually. Case fatality rates often exceeded 30 percent, and no effective treatment existed. Parents would watch their children suffer and die, helpless against this relentless disease.

Behring's path to the diphtheria cure began when he joined Robert Koch's laboratory at the Institute for Infectious Diseases in Berlin in 1890. Koch, the discoverer of the tuberculosis bacillus, was the dominant figure in bacteriology, and his laboratory was the epicenter of microbiological research worldwide. It was here that Behring, working alongside Kitasato Shibasaburo, made his first landmark discovery: that blood serum from animals that had been immunized against tetanus or diphtheria could confer passive immunity to other animals.

\section{The Discovery}

The intellectual foundation for Behring's work rested on a growing body of evidence that the body possessed natural defense mechanisms against infectious organisms\index{Passive immunity|textbf}. Pasteur and others had demonstrated that attenuated or killed microorganisms could immunize animals against subsequent infection. But the mechanism of this immunity was poorly understood. Behring's key insight was that the protective agent was not the microorganism itself but rather a substance produced in the blood in response to the infection.

Working initially with tetanus, Behring and Kitasato demonstrated in 1890 that serum from rabbits that had been immunized with sublethal doses of tetanus toxin could protect other rabbits from otherwise fatal doses of the same toxin. The protective substance was located in the serum, the liquid fraction of blood, and was distinct from the blood cells themselves. Behring called this protective substance \emph{Antitoxin}, coining a term that would become central to immunology. He proposed that toxins produced by bacteria could be neutralized by antitoxins produced by the body in response to immunization.

The translation of this laboratory finding into a practical therapy for diphtheria required years of painstaking work\index{Antitoxin|textbf}. Behring collaborated with the pharmaceutical company Farbwerke Hoechst to develop a process for producing antitoxin on a large scale. Horses were immunized with gradually increasing doses of diphtheria toxin over a period of months. The horses tolerated the immunization well, and their blood serum developed high concentrations of antitoxin. The serum was then collected, purified, and tested for safety and potency.

The first clinical trials of diphtheria antitoxin were conducted in 1891 at the Berlin Charity Hospital, under the direction of the pediatrician Clemens Freiherr von Pirquet (who would later make his own important contributions to immunology). The results were dramatic\index{Clinical trials|textbf}. Children who received the antitoxin showed rapid improvement, and mortality rates dropped significantly compared to untreated patients. The antitoxin did not kill the diphtheria bacillus itself; rather, it neutralized the deadly toxin that the bacteria produced, allowing the body's own defenses to clear the infection.

Behring's antitoxin represented the first truly effective therapy for a major infectious disease. Its impact was immediate and significant. By the mid-1890s, diphtheria antitoxin was being produced and distributed across Europe and America. Mortality from diphtheria fell dramatically---in many cities, death rates dropped from 30--40 percent to under 10 percent within a few years of the antitoxin's introduction. The success of diphtheria antitoxin also validated the broader concept of serum therapy, inspiring efforts to develop antitoxins for other diseases including tetanus, scarlet fever, and pneumonia.

The socioeconomic impact of diphtheria antitoxin was profound. Hospitals, which had previously been places where families gathered to watch loved ones die, became centers of hope and recovery. The reduction in child mortality allowed families to retain their children, strengthening household economies and reducing the burden on orphanages and charitable institutions. Governments recognized the public health value of antitoxin and began to subsidize its production and distribution, establishing the principle that the state had a responsibility to provide life-saving treatments\index{Public health policy|textbf}.

Beyond the antitoxin itself, Behring's work established fundamental principles of immunology that remain valid today\index{Immunology|textbf}. His demonstration that immunity could be passively transferred through serum laid the groundwork for the development of monoclonal antibodies, intravenous immunoglobulin therapy, and many other immunological interventions. His concept of the antitoxin---a molecule that specifically neutralizes a toxin---was the intellectual precursor to modern antibody-based therapeutics.

In recognition of this transformative work, Emil von Behring was awarded the very first Nobel Prize in Physiology or Medicine in 1901, making him the inaugural laureate in what would become the most prestigious award in medicine\index{Nobel Prize|textbf}. The award ceremony took place at the Karolinska Institutet in Stockholm on December 10, 1901. Behring was recognized not only for his specific discovery but for founding an entirely new field of medical science: immunology. The ceremony was attended by King Oscar II of Sweden and Norway, and Behring's acceptance speech emphasized that the prize belonged not to him alone but to the entire scientific community that had made the discovery possible.

The Nobel Committee's deliberations had been extensive. Robert Koch, Behring's former mentor, was also considered for the prize, but the committee ultimately decided that Behring's work on serum therapy was more directly applicable to clinical medicine and had saved more lives in the short term\index{Nobel Committee|textbf}. The inaugural ceremony established traditions that continue today: the Nobel Lecture, the formal banquet at Stockholm City Hall, and the awarding of a gold medal and diploma. Behring was reportedly surprised by the honor, having expected that Koch would receive the prize, and his modesty in accepting it endeared him to the Swedish royal family.

\section{Impact on Modern Medicine}

The direct descendants of Behring's antitoxin therapy are still in use today, more than a century later\index{Monoclonal antibodies|textbf}. Diphtheria antitoxin remains the specific treatment for suspected diphtheria and is included in the essential medicines list of the World Health Organization. While diphtheria has been largely controlled in developed nations through vaccination with the diphtheria toxoid (part of the DTaP and Tdap combination vaccines), the antitoxin is still manufactured and stockpiled as an emergency treatment.

The broader principle of passive immunization that Behring pioneered---transferring pre-formed antibodies from one individual to another---has become a cornerstone of modern medicine\index{Antivenom|textbf}. Antivenoms for snake and spider bites are produced by the same principle: animals are immunized with venom, and their serum, containing neutralizing antibodies, is used to treat envenomation in humans. Intravenous immunoglobulin (IVIG), derived from pooled human plasma, is used to treat a wide range of conditions including primary immunodeficiencies, autoimmune diseases, and certain infections\index{IVIG|textbf}.

More recently, monoclonal antibody therapies---engineered proteins that mimic the specificity of natural antibodies---have revolutionized the treatment of cancer, autoimmune disorders, and infectious diseases including COVID-19\index{COVID-19|textbf}. In 2023, over 100 monoclonal antibodies were approved for clinical use worldwide. The COVID-19 pandemic saw the rapid development of monoclonal antibody therapies (e.g., bamlanivimab, casirivimab) that neutralize the SARS-CoV-2 virus.

The scale of antivenom production illustrates the breadth of Behring's influence. Approximately 5.4 million snakebite envenomations occur annually worldwide, causing 81,000--138,000 deaths\index{Snakebite|textbf}. Modern antivenoms, produced using the same principles Behring established, are manufactured in horses and sheep, with serum purified and concentrated for therapeutic use. Self-administered subcutaneous immunoglobulin (SCIG) therapy, developed from Behring's passive immunization principle, allows patients with immunodeficiencies to receive treatment at home, dramatically improving quality of life.

Behring's work also catalyzed the development of the modern pharmaceutical industry\index{Pharmaceutical industry|textbf}. The large-scale production of diphtheria antitoxin required standardized manufacturing processes, quality control, and potency testing, setting precedents for the regulation and production of biological medicines that persist to this day. The development of antitoxin production facilities in Germany, France, and the United States marked the beginning of the biologics industry, which now encompasses vaccines, blood products, gene therapies, and recombinant proteins.

The success of antitoxin therapy also stimulated research into the nature of the immune response itself. Scientists asked: How does the body produce antitoxins? What is the chemical nature of these protective substances? These questions drove the development of biochemistry and immunochemistry in the twentieth century, leading eventually to the discovery of the structure and function of antibodies by Rodney Porter and Gerald Edelman in the 1950s and 1960s, and ultimately to the era of engineered monoclonal antibodies that dominates much of modern therapeutics\index{Rodney Porter|textbf}.

The regulatory framework for biological medicines also traces its origins to Behring's work. The Biologics Control Act of 1902, passed in the United States in response to tragic incidents involving contaminated antitoxins, established the requirement that biological products be pure, safe, and potent\index{Biologics Control Act|textbf}. This legislation laid the foundation for modern regulatory oversight of vaccines, blood products, and gene therapies, and the principles it established continue to govern the approval of new biological medicines today.

\section{Legacy}

Emil von Behring's legacy is multifaceted\index{Behring, Emil von|textbf}. His discovery of the antitoxin principle saved countless lives and inaugurated the era of immunological therapy. The Nobel Prize in Physiology or Medicine that he received in 1901 was not merely a personal honor but a recognition that immunology had become one of the major frontiers of medical science.

Behring went on to make further contributions to medicine, including work on the diphtheria toxoid vaccine---the attenuated toxin that, when injected, stimulated active immunity against diphtheria without causing the disease. This toxoid vaccine, developed in collaboration with Gaston Ramon, became the basis for the diphtheria component of the universal childhood vaccination programs that have virtually eliminated the disease in much of the world. Behring also established a pharmaceutical research institute in Marburg, Germany, which became a center for immunological research.

The institutions and concepts that Behring helped create endure. The Paul Ehrlich Institute, which succeeded the institute that Behring co-founded, remains one of the world's leading centers for research on vaccines, blood products, and immunological therapies\index{Paul Ehrlich Institute|textbf}. The concept of the antitoxin, and the broader understanding of passive immunity that it embodied, continue to inform the design of new immunotherapies. In an era of monoclonal antibodies, checkpoint inhibitors, and CAR-T cell therapy, the intellectual lineage can be traced back to that Berlin laboratory in the 1890s where Emil von Behring first demonstrated that serum could cure disease.

Behring's story also illustrates a recurring theme in the history of medicine: that therapeutic breakthroughs often arise from fundamental research conducted without immediate clinical application in mind\index{Translational research|textbf}. Behring did not set out to develop a cure for diphtheria; he was studying the basic biology of immunity. But by following the thread of scientific inquiry wherever it led, he arrived at a discovery that transformed the treatment of one of the era's most devastating diseases. In this respect, Behring exemplifies the principle that investment in basic science is ultimately the most reliable path to medical progress.

Behring's relationship with Kaiser Wilhelm II was close, and the Kaiser granted him the hereditary title of von Behring in 1901, elevating him to the nobility\index{Kaiser Wilhelm II|textbf}. This honor reflected the esteem in which Behring was held not only by the scientific community but by the German state. Beyond diphtheria, Behring continued to work on plague antitoxin and the phenomenon of anaphylaxis---the severe allergic reaction to foreign proteins---demonstrating his ongoing commitment to understanding and combating infectious diseases.

\section{Modern Developments}

Behring's antitoxin principle has evolved into the era of monoclonal antibodies---engineered proteins that target specific molecules with unprecedented precision\index{Therapeutics|textbf}. In 2023, over 100 monoclonal antibodies were approved for clinical use, treating cancers, autoimmune diseases, and infectious diseases. CAR-T cell therapy, which engineers a patient's own immune cells to attack cancer, represents the latest evolution of Behring's insight that immunity can be harnessed therapeutically. The principle of passive immunity remains central to modern medicine, from antivenoms to RSV prophylaxis (nirsevimab).


\section{Bibliography}

\begin{thebibliography}{9}
\bibitem{nobel-1901}
Nobel Prize Outreach, ``The Nobel Prize in Physiology or Medicine 1901,''\emph{NobelPrize.org}.\\
\url{https://www.nobelprize.org/prizes/medicine/1901/summary/}

\bibitem{lecture-1901}
Emil von Behring, ``Nobel Lecture,''\emph{NobelPrize.org}. Winner-authored primary source; downloadable from the official Nobel Prize archive.\\
\url{https://www.nobelprize.org/prizes/medicine/1901/behring/lecture/}
\end{thebibliography}
